%! Random Sampling Methods — Unit 3, Lesson 3.2 — Table Activity: Which Plan Fits? — STUDENT SHEETS
% ONE-OFF, STANDALONE HANDOUT. Not a packet component: `table_activity` is deliberately NOT a
% name in shared/lesson.mk's STUDENT_ORDER, so `make` never builds it and no packet ever
% merges it. The teacher prints it only on a day it is used (see the optional box at the end
% of ../main.tex). Build it by hand, out of tree:
%   cd unit03/lesson02/table_activity && TEXINPUTS="$(git rev-parse --show-toplevel)/shared//:" \
%     latexmk -xelatex -outdir=../../../target/unit03/lesson02/table_activity main.tex
%
% THREE TABLES, THREE CARDS, ONE BEST METHOD EACH — each sheet is exactly two pages
% (print duplex, one sheet per student at that table):
%   Table 1  Pine Ridge County smoke detectors  -> CLUSTER   (no list, in-person visits,
%            map sections each a mix like the county)
%   Table 2  StreamLine music app ratings       -> SRS       (full numbered list, free to reach
%            anyone, nothing known to stratify on)
%   Table 3  Northfield Foods commute distance  -> STRATIFIED by work site (commutes alike
%            inside a site) — and THE CRUX: "pick one site as a cluster"
% Contexts are new: not Lakeview / Riverside (notes), the hospital (homework), or farms /
% textbooks / the airplane (AP Practice).
%
% PAGE LOCKSTEP with ../table_activity_key/: this file is the key with every \answerspace,
% \cellans and \fillin second argument emptied (the key fills them) — nothing else differs but the style package
% and the header line.
\documentclass[12pt]{article}
\usepackage{apstats-article}
\usepackage{apstats-boxes}
\newcommand{\answerspace}[2]{\par\nopagebreak\noindent\begin{minipage}[t][#1][t]{\linewidth}%
  \color{keyred}\bfseries\footnotesize #2\end{minipage}\par}
\newcommand{\cellans}[2]{\begin{minipage}[t][#1][t]{\linewidth}\raggedright
  \color{keyred}\bfseries\footnotesize #2\end{minipage}}
\newcommand{\fillin}[2]{\underline{\makebox[#1][l]{\ \smash{\color{keyred}\bfseries #2}}}}
\newcolumntype{L}[1]{>{\raggedright\arraybackslash}p{#1}}

\begin{document}

% ════════════════════════════════════════════════════════════════════════════
% TABLE 1 — PINE RIDGE COUNTY — CLUSTER
% ════════════════════════════════════════════════════════════════════════════
\pageheader{Unit 3, Lesson 3.2}{Table Activity: Which Plan Fits? \quad --- \quad Table 1}

\boxguard[10]
\begin{headlinebox}{sky}
\small
\textbf{Your table is the survey team.} Read your card, try all three sampling methods on it,
and choose the one that fits. Be ready to defend your choice to the class in one minute ---
\emph{in context}, using the words \emph{alike inside}, \emph{looks like the whole}, or
\emph{cost}.
\end{headlinebox}

\vspace{0.08in}

\boxguard[12]
\begin{scenariobox}[Table 1 card: Smoke Detectors in Pine Ridge County]{navy}
\small
\textbf{Statistical question:} What proportion of homes in Pine Ridge County have a working
smoke detector?
\begin{itemize}[leftmargin=1.3em, itemsep=2pt, topsep=4pt]
  \item About 9{,}000 homes are spread over 600 square miles of farmland, hills, and small
        towns.
  \item \textbf{There is no list of homes.} The county map does divide the county into 150
        \emph{map sections} of about 60 homes each.
  \item Every map section is a mix: old and new houses, owners and renters, families of every
        income.
  \item A fire inspector must visit each home \emph{in person} and press the detector's test
        button. Driving from one map section to another can take an hour.
  \item The fire department can afford to check about \textbf{300 homes}.
\end{itemize}
\end{scenariobox}

\vspace{0.08in}

\boxguard[10]
\begin{scenariobox}[1. Read the population]{navy}
\small
\begin{enumerate}[label=\textbf{\alph*.}, leftmargin=1.6em, itemsep=4pt]
  \item What are the individuals? What is the variable?
        \answerspace{1.1cm}{}
  \item Is there a list of every individual? What does it cost to reach one?
        \answerspace{1.3cm}{}
  \item Are there groups that are \emph{alike inside} on the response? Groups that each
        \emph{look like the whole} population?
        \answerspace{1.3cm}{}
\end{enumerate}
\end{scenariobox}

\clearpage

\boxguard[16]
\begin{scenariobox}[2. Try all three methods on your card]{navy}
\small
\renewcommand{\arraystretch}{1.2}
\begin{tabular}{@{}L{2.1cm}|L{6.1cm}|L{6.1cm}@{}}
\textbf{Method} & \textbf{How would it work here?} & \textbf{Does it fit? Why or why not?} \\
\hline
\textbf{Simple random sample} &
\cellans{2.0cm}{} &
\cellans{2.0cm}{} \\
\hline
\textbf{Stratified random sample} &
\cellans{2.0cm}{} &
\cellans{2.0cm}{} \\
\hline
\textbf{Cluster sample} &
\cellans{2.0cm}{} &
\cellans{2.0cm}{} \\
\end{tabular}
\end{scenariobox}

\vspace{0.08in}

\boxguard[12]
\begin{scenariobox}[3. Our choice]{navy}
\small
\begin{enumerate}[label=\textbf{\alph*.}, leftmargin=1.6em, itemsep=4pt]
  \item The best method for our card: \fillin{8.5cm}{}
  \item Defend it \emph{in context}: why does this method fit \emph{this} population better than
        the other two?
        \answerspace{1.8cm}{}
  \item Describe \emph{exactly} how to choose the sample, step by step.
        \answerspace{1.8cm}{}
\end{enumerate}
\end{scenariobox}

\vspace{0.08in}

\boxguard[8]
\begin{scenariobox}[4. Listen to the other tables]{navy}
\small
\renewcommand{\arraystretch}{1.2}
\begin{tabular}{@{}L{1.3cm}|L{4.0cm}|L{2.9cm}|L{5.8cm}@{}}
\textbf{Table} & \textbf{Their question} & \textbf{Best method} & \textbf{The reason, in one phrase} \\
\hline
\textbf{2} & \cellans{1.2cm}{} & \cellans{1.2cm}{} &
\cellans{1.2cm}{} \\
\hline
\textbf{3} & \cellans{1.2cm}{} & \cellans{1.2cm}{} &
\cellans{1.2cm}{} \\
\end{tabular}

\vspace{6pt}
\centering\bfseries
Strata are alike inside: sample from \fillin{1.4cm}{} one.\par
Clusters each look like the whole: take \fillin{1.0cm}{} of \fillin{1.3cm}{} of them.
\end{scenariobox}

% ════════════════════════════════════════════════════════════════════════════
% TABLE 2 — STREAMLINE — SIMPLE RANDOM SAMPLE
% ════════════════════════════════════════════════════════════════════════════
\clearpage
\pageheader{Unit 3, Lesson 3.2}{Table Activity: Which Plan Fits? \quad --- \quad Table 2}

\boxguard[10]
\begin{headlinebox}{sky}
\small
\textbf{Your table is the survey team.} Read your card, try all three sampling methods on it,
and choose the one that fits. Be ready to defend your choice to the class in one minute ---
\emph{in context}, using the words \emph{alike inside}, \emph{looks like the whole}, or
\emph{cost}.
\end{headlinebox}

\vspace{0.08in}

\boxguard[12]
\begin{scenariobox}[Table 2 card: Rating a New Feature on StreamLine]{navy}
\small
\textbf{Statistical question:} What is the mean rating, on a scale of 1 to 10, that StreamLine
subscribers give the app's new \emph{Daily Mix} feature?
\begin{itemize}[leftmargin=1.3em, itemsep=2pt, topsep=4pt]
  \item StreamLine, a music app, has 50{,}000 subscribers all over the country.
  \item \textbf{The company has a list of every subscriber}, numbered 00001 to 50000. The list
        records nothing else about them --- no age, no location, no listening habits.
  \item A survey can be sent inside the app to any subscriber \emph{at no cost}, and it reaches
        them wherever they live.
  \item The company wants \textbf{200 ratings}.
\end{itemize}
\end{scenariobox}

\vspace{0.08in}

\boxguard[10]
\begin{scenariobox}[1. Read the population]{navy}
\small
\begin{enumerate}[label=\textbf{\alph*.}, leftmargin=1.6em, itemsep=4pt]
  \item What are the individuals? What is the variable?
        \answerspace{1.1cm}{}
  \item Is there a list of every individual? What does it cost to reach one?
        \answerspace{1.3cm}{}
  \item Are there groups that are \emph{alike inside} on the response? Groups that each
        \emph{look like the whole} population?
        \answerspace{1.3cm}{}
\end{enumerate}
\end{scenariobox}

\clearpage

\boxguard[16]
\begin{scenariobox}[2. Try all three methods on your card]{navy}
\small
\renewcommand{\arraystretch}{1.2}
\begin{tabular}{@{}L{2.1cm}|L{6.1cm}|L{6.1cm}@{}}
\textbf{Method} & \textbf{How would it work here?} & \textbf{Does it fit? Why or why not?} \\
\hline
\textbf{Simple random sample} &
\cellans{2.0cm}{} &
\cellans{2.0cm}{} \\
\hline
\textbf{Stratified random sample} &
\cellans{2.0cm}{} &
\cellans{2.0cm}{} \\
\hline
\textbf{Cluster sample} &
\cellans{2.0cm}{} &
\cellans{2.0cm}{} \\
\end{tabular}
\end{scenariobox}

\vspace{0.08in}

\boxguard[12]
\begin{scenariobox}[3. Our choice]{navy}
\small
\begin{enumerate}[label=\textbf{\alph*.}, leftmargin=1.6em, itemsep=4pt]
  \item The best method for our card: \fillin{8.5cm}{}
  \item Defend it \emph{in context}: why does this method fit \emph{this} population better than
        the other two?
        \answerspace{1.8cm}{}
  \item Describe \emph{exactly} how to choose the sample, step by step.
        \answerspace{1.8cm}{}
\end{enumerate}
\end{scenariobox}

\vspace{0.08in}

\boxguard[8]
\begin{scenariobox}[4. Listen to the other tables]{navy}
\small
\renewcommand{\arraystretch}{1.2}
\begin{tabular}{@{}L{1.3cm}|L{4.0cm}|L{2.9cm}|L{5.8cm}@{}}
\textbf{Table} & \textbf{Their question} & \textbf{Best method} & \textbf{The reason, in one phrase} \\
\hline
\textbf{1} & \cellans{1.2cm}{} & \cellans{1.2cm}{} &
\cellans{1.2cm}{} \\
\hline
\textbf{3} & \cellans{1.2cm}{} & \cellans{1.2cm}{} &
\cellans{1.2cm}{} \\
\end{tabular}

\vspace{6pt}
\centering\bfseries
Strata are alike inside: sample from \fillin{1.4cm}{} one.\par
Clusters each look like the whole: take \fillin{1.0cm}{} of \fillin{1.3cm}{} of them.
\end{scenariobox}

% ════════════════════════════════════════════════════════════════════════════
% TABLE 3 — NORTHFIELD FOODS — STRATIFIED (THE CRUX: a site is not a cluster)
% ════════════════════════════════════════════════════════════════════════════
\clearpage
\pageheader{Unit 3, Lesson 3.2}{Table Activity: Which Plan Fits? \quad --- \quad Table 3}

\boxguard[10]
\begin{headlinebox}{sky}
\small
\textbf{Your table is the survey team.} Read your card, try all three sampling methods on it,
and choose the one that fits. Be ready to defend your choice to the class in one minute ---
\emph{in context}, using the words \emph{alike inside}, \emph{looks like the whole}, or
\emph{cost}.
\end{headlinebox}

\vspace{0.08in}

\boxguard[12]
\begin{scenariobox}[Table 3 card: Commuting to Northfield Foods]{navy}
\small
\textbf{Statistical question:} What is the mean one-way commute distance, in miles, for
Northfield Foods employees?
\par\vspace{4pt}
{\centering\footnotesize
\renewcommand{\arraystretch}{1.2}
\begin{tabular}{@{}l c l@{}}
\rowcolor{sky}
\textbf{Work site} & \textbf{Employees} & \textbf{Where most of them live} \\
\hline
Downtown plant & 1{,}200 & in the city, close to the plant \\
Riverbend farm, 40 miles out of the city & 400 & in the city --- they drive out \\
Lakeside office, in the suburbs & 200 & in the nearby suburbs \\
\end{tabular}\par}
\begin{itemize}[leftmargin=1.3em, itemsep=2pt, topsep=4pt]
  \item \textbf{Human Resources has a list of all 1{,}800 employees} and the site where each
        one works.
  \item Any employee can be surveyed by company email \emph{at no cost}.
  \item The company wants \textbf{90 responses}.
\end{itemize}
\end{scenariobox}

\vspace{0.08in}

\boxguard[10]
\begin{scenariobox}[1. Read the population]{navy}
\small
\begin{enumerate}[label=\textbf{\alph*.}, leftmargin=1.6em, itemsep=4pt]
  \item What are the individuals? What is the variable?
        \answerspace{1.1cm}{}
  \item Is there a list of every individual? What does it cost to reach one?
        \answerspace{1.3cm}{}
  \item Are there groups that are \emph{alike inside} on the response? Groups that each
        \emph{look like the whole} population?
        \answerspace{1.3cm}{}
\end{enumerate}
\end{scenariobox}

\clearpage

\boxguard[16]
\begin{scenariobox}[2. Try all three methods on your card]{navy}
\small
\renewcommand{\arraystretch}{1.2}
\begin{tabular}{@{}L{2.1cm}|L{6.1cm}|L{6.1cm}@{}}
\textbf{Method} & \textbf{How would it work here?} & \textbf{Does it fit? Why or why not?} \\
\hline
\textbf{Simple random sample} &
\cellans{2.0cm}{} &
\cellans{2.0cm}{} \\
\hline
\textbf{Stratified random sample} &
\cellans{2.0cm}{} &
\cellans{2.0cm}{} \\
\hline
\textbf{Cluster sample} &
\cellans{2.0cm}{} &
\cellans{2.0cm}{} \\
\end{tabular}
\end{scenariobox}

\vspace{0.08in}

\boxguard[12]
\begin{scenariobox}[3. Our choice]{navy}
\small
\begin{enumerate}[label=\textbf{\alph*.}, leftmargin=1.6em, itemsep=4pt]
  \item The best method for our card: \fillin{8.5cm}{}
  \item Defend it \emph{in context}: why does this method fit \emph{this} population better than
        the other two?
        \answerspace{1.8cm}{}
  \item Describe \emph{exactly} how to choose the sample, step by step.
        \answerspace{1.8cm}{}
\end{enumerate}
\end{scenariobox}

\vspace{0.08in}

\boxguard[8]
\begin{scenariobox}[4. Listen to the other tables]{navy}
\small
\renewcommand{\arraystretch}{1.2}
\begin{tabular}{@{}L{1.3cm}|L{4.0cm}|L{2.9cm}|L{5.8cm}@{}}
\textbf{Table} & \textbf{Their question} & \textbf{Best method} & \textbf{The reason, in one phrase} \\
\hline
\textbf{1} & \cellans{1.2cm}{} & \cellans{1.2cm}{} &
\cellans{1.2cm}{} \\
\hline
\textbf{2} & \cellans{1.2cm}{} & \cellans{1.2cm}{} &
\cellans{1.2cm}{} \\
\end{tabular}

\vspace{6pt}
\centering\bfseries
Strata are alike inside: sample from \fillin{1.4cm}{} one.\par
Clusters each look like the whole: take \fillin{1.0cm}{} of \fillin{1.3cm}{} of them.
\end{scenariobox}

\end{document}
